La modernización de electrónica envejecida y crítica para la seguridad suele ser más rentable y más sostenible desde el punto de vista medioambiental que su sustitución completa, siempre que la equivalencia funcional con el sistema original pueda demostrarse y el sistema modernizado pueda volver a certificarse conforme a normas actuales y más exigentes. Este estudio de caso, que da continuidad a trabajos previos, aborda este reto tomando como ejemplo un módulo de prueba y diagnóstico heredado, perteneciente a un sistema antideslizamiento y utilizado en el ferrocarril urbano europeo.

Una placa controladora de 8~bits heredada, construida en torno al Intel~MCS-48, se modernizó con una matriz de puertas programable en campo (FPGA) \texttt{Actel~A3P1000} basada en memoria flash, que emula la unidad central de procesamiento (CPU) y la lógica de entrada/salida (E/S) originales, complementada con un microcontrolador (µC) Arm\textsuperscript{\textregistered} de 32~bits aislado galvánicamente, un \texttt{STM32F401RET6}, para un registro de datos ampliado. Un objetivo central era desarrollar un marco metodológico de pruebas adecuado para una aplicación más amplia en la modernización de sistemas críticos para la seguridad. Su núcleo es una estrategia de verificación en dos niveles derivada del modelo en~V: (1) Pruebas de componentes e interfaces: la simulación determinista en lenguaje de descripción de hardware (HDL), las comprobaciones JTAG y las mediciones de temporización en el hardware examinan la corrección lógica, la integridad básica del hardware y la temporización de las interfaces (por ejemplo, perro guardián y barrido de frecuencia). (2) Pruebas a nivel de sistema: las pruebas basadas en escenarios, realizadas en un banco de pruebas específico~\cite{dolata2023} con casos de prueba definidos en YAML, examinan el comportamiento integrado (memoria de fallos, autopruebas, diagnóstico). La simulación y las pruebas integradas basadas en escenarios indican la equivalencia funcional del tratamiento de fallos del núcleo crítico para la seguridad (detección, almacenamiento, visualización y lectura de fallos), pero la única medición comparativa con la placa original muestra una desviación del barrido de frecuencia, que en la emulación termina unos 70~Hz más abajo (unos 850~Hz frente a unos 920~Hz) y transcurre más rápido que la variación de frecuencia del original, de unos 180~ms; esta desviación, que sigue abierta, requiere una nueva medición de la placa original o una comprobación de su circuito de barrido, y la evidencia aún no está completa: las simulaciones y mediciones archivadas datan de 2023 y son, por tanto, anteriores a la reconstrucción; en la unidad reconstruida, un equipo que continuó los trabajos ejecutó el conjunto de pruebas del banco sin registro archivado, y varios escenarios, entre ellos la autoprueba, no están documentados para esa unidad. Seis defectos de integración documentados durante la reconstrucción se clasifican según cómo se encontraron; no salieron a la luz en una aplicación planificada del método en dos niveles, y ningún método por sí solo los encontró todos. La metodología se puso en correspondencia con las normas IEC~61508~Ed.~2, EN~50129:2018, EN~50155:2021, EN~50716:2023 e IEEE~1012:2024. Aunque las instalaciones completas de pruebas de hardware en el bucle (HIL) quedaban fuera del alcance de este proyecto, el enfoque sigue los principios HIL para facilitar una integración futura.

El reacondicionamiento de la electrónica existente del material rodante puede reducir costes y evitar las emisiones incorporadas que conlleva la fabricación de componentes nuevos, y las modernizaciones a nivel de placa obedecen a la misma lógica. Una función adaptativa de protección antideslizamiento de las ruedas (WSP), destinada también a flotas existentes y certificada conforme a la norma vigente EN~15595, constituye un precedente; la validación de la metodología en otros casos de modernización queda como trabajo futuro.
